\documentclass[10pt]{article}
\usepackage[margin=0.9in]{geometry}
\usepackage{graphicx}
\usepackage{booktabs}
\usepackage{amsmath}
\usepackage{caption}
\usepackage[hidelinks]{hyperref}
\usepackage{float}
\captionsetup{font=small}
\setlength{\parskip}{4pt}
\setlength{\parindent}{0pt}

\title{\vspace{-2.5em}Deep Snake on FLAME-1: Detector-Coupled vs.\ Detector-Free\\
\large Re-running the deep-contour methods on a visible-flame dataset}
\author{CS269 Project --- Deformable Models}
\date{\today}

\begin{document}
\maketitle
\vspace{-1.5em}

\section{Summary}

We re-ran the two deep-contour ``snake'' methods on \textbf{FLAME-1} (Sycan Marsh,
2{,}001 RGB frames with hand-labeled \emph{visible-flame} PNG masks; contiguous
$70/15/15$ split $\to$ 1{,}400 train / 300 val / 301 test). Two implementations
were compared:

\begin{itemize}
\item \texttt{deep\_snake\_simple} --- the original detector-free design: a shared
  U-Net produces a coarse mask whose connected components seed the contour, which
  is then refined by circular-convolution snake heads.
\item \texttt{deep\_snake\_paper} --- a paper-faithful re-implementation that
  \emph{decouples detection from contour refinement}: a separately-trained MMDet
  \textbf{CenterNet} detector localizes fire boxes, an octagon is initialized from
  each box, and the snake deforms it. The snake is trained on ground-truth boxes;
  at test time it consumes the detector's predicted boxes.
\end{itemize}

\begin{table}[H]
\centering
\caption{Test results on FLAME-1 (301 frames, IoU/Dice on binary masks).
Detector trained 10 epochs; snakes 30 epochs.}
\begin{tabular}{lccc}
\toprule
Method & IoU & Dice & Frames with a prediction \\
\midrule
\texttt{deep\_snake\_simple} (detector-free) & \textbf{0.658} & \textbf{0.792} & 301 / 301 (100\%) \\
\texttt{deep\_snake\_paper} (CenterNet $\to$ snake) & 0.050 & 0.080 & 73 / 301 (24\%) \\
\bottomrule
\end{tabular}
\end{table}

\section{Headline finding: the task, not the architecture}

The single most important result is the contrast with the recorded \textbf{FLAME-3}
numbers. On FLAME-3 (where ground truth is the paired thermal frame thresholded at
$\ge150\,^\circ$C---a smoke-occluded hot region that is essentially \emph{invisible}
in RGB), the deep snake scored IoU $0.032$, below even the learned U-Net ($0.147$)
and DALS ($0.122$), and far below oracle-initialized GAC ($0.365$). The color floor
was $0.004$.

On FLAME-1, where the label \emph{is} the orange flame visible in RGB,
\texttt{deep\_snake\_simple} jumps to \textbf{IoU 0.658 / Dice 0.792}---a
$\sim$20$\times$ improvement, and higher than \emph{any} FLAME-3 method including the
oracle-seeded classical contours, with \emph{no} oracle prior. This is direct
evidence that Deep Snake's weak FLAME-3 showing was a property of the
\emph{task} (RGB cannot see thermal-defined fire), not of the method.

\begin{table}[H]
\centering
\caption{FLAME-3 (recorded, RESULTS.md) vs.\ FLAME-1 (this run).}
\begin{tabular}{lcc}
\toprule
 & FLAME-3 IoU & FLAME-1 IoU \\
\midrule
Color floor (R$-$G)            & 0.004 & --- \\
U-Net (RGB)                    & 0.147 & --- \\
DALS (RGB)                     & 0.122 & --- \\
Oracle-init GAC                & 0.365 & --- \\
Deep Snake (detector-free)     & 0.032 & \textbf{0.658} \\
\bottomrule
\end{tabular}
\end{table}

\section{Why the paper-faithful pipeline underperforms here}

\texttt{deep\_snake\_paper} scored only $0.050$ overall---but this is a
\textbf{detector-recall bottleneck}, not a snake failure:

\begin{itemize}
\item The CenterNet detector reached \textbf{mAP $0.120$ / mAP@50 $0.430$} after 10
  epochs. Notably, this is a \emph{working} detector: on FLAME-3 the same pipeline
  trained to mAP $0.000$ (it never learned), and was unusable. The fix that made it
  train at all was correcting the learning rate (the base config's LR of $0.01$ is
  calibrated for an effective batch of 128; at our batch of 4 the heatmap head
  collapsed---auto-scaling to $\approx 3\times10^{-4}$ resolved it).
\item At test time the detector fires on only \textbf{24\% of frames}. Where it
  misses, the snake receives no box and produces an empty mask, dragging the mean
  to $0.050$.
\item \emph{Conditioned on the detector firing}, the snake reaches IoU
  \textbf{0.206}---over 4$\times$ the unconditional score. Its oracle-box validation
  IoU was $0.180$. So the contour refinement works; the pipeline is gated by how
  often the detector finds the fire.
\end{itemize}

The decoupled design pays for an explicit detection stage that, at 10 epochs on
small/fragmented fire boxes (median $\sim$2 dominant components per frame), has low
recall. The detector-free variant sidesteps this by seeding from a dense coarse
mask, which on a visible-flame dataset is far more reliable.

\section{Takeaways}

\begin{enumerate}
\item \textbf{Deep Snake is competitive when the target is visible in RGB.} The
  detector-free snake hits IoU 0.658 on FLAME-1 vs.\ 0.032 on FLAME-3---the dataset,
  not the model, drove the original poor result.
\item \textbf{The paper-faithful detector$\to$snake pipeline is recall-limited.} The
  snake itself refines well (0.206 IoU when given a box, 0.180 oracle), but a
  10-epoch CenterNet only localizes 24\% of test frames. More detector training
  (epochs were capped for runtime; it was still improving) and less-fragmented box
  targets are the obvious levers.
\item \textbf{The pipeline is now functional end-to-end}, which it was not on
  FLAME-3 (detector mAP 0). The remaining gap is detector capacity/training, a
  tractable next step.
\end{enumerate}

\vspace{0.5em}
\footnotesize
Artifacts: \texttt{results/rerun\_flame1\_20260601\_193349/} (per-frame CSVs, logs,
detector checkpoint, provenance). Detector: 10 epochs; snakes: 30 epochs; RTX 4090.
Recorded additively; committed FLAME-3 results and \texttt{RESULTS.md} unchanged.

\end{document}
